% !TeX program = xelatex
% UTF-8; standalone XeLaTeX/XePersian transcription.
% Source: Google Drive / Notebook / pertussis, in filename/page order.
% images/pic01.jpg = 20261018_193949333.jpg (1ulbEfK2B2aEZNIPAsO0mCwrEkMEeH7qE)
% images/pic02.jpg = 20261018_193953633.jpg (1bCdL0lNql2qLXtgcptvBnGwfOFZtctEi)
% Source wording, abbreviations, numbers, corrections and deleted text are retained.
% Arrows are represented by lists or the relation command below.
% The source's abbreviated thousands in the WBC range are expanded to full digits.
\documentclass[11pt,a4paper]{article}
\usepackage[margin=18mm]{geometry}
\usepackage{amsmath,amssymb}
\usepackage{enumitem}
\usepackage{xcolor}
\usepackage[normalem]{ulem}
\usepackage{xepersian}
\settextfont{Amiri}
\setlatintextfont{DejaVu Serif}
\setlength{\parindent}{0pt}
\setlength{\parskip}{4pt}
\setlength{\emergencystretch}{3em}
\setlist[itemize]{itemsep=3pt,topsep=3pt,leftmargin=*}
\newcommand{\en}[1]{\lr{#1}}
\newcommand{\hl}[1]{\begingroup\setlength{\fboxsep}{1pt}\colorbox{yellow!40}{\strut #1}\endgroup}
\newcommand{\rel}{\enspace\lr{\(\longleftarrow\)}\enspace}
\newcommand{\sheet}[1]{%
  \section*{تصویر #1: ۱۰۸ ـ پرتوسیس}
  {\small شمارهٔ دست‌نویس پایین برگه: #1؛ نشان چاپی: \en{GNOTE}.\par}}

\begin{document}

\sheet{۱}
{\small نوشته‌های چاپی بالای برگه: \en{Date:}\quad\en{Subject:}\par}

\begin{itemize}
\item سیاه‌سرفه (\en{Whooping cough}) \rel بوردتلا پرتوسیس.
\begin{itemize}
\item دورهٔ انکوبیشن \rel \hl{۷ تا ۱۰ روز}.
\item بیشترین احتمال انتقال \rel \hl{۲ هفتهٔ اول سرفه‌ها}.
\item سن شایع \rel زیر ۶ ماه که هنوز واکسن نزده‌اند.
\end{itemize}
\end{itemize}

\subsection*{کلینیک تیپیک}
\hl{کودک ۱ تا ۱۰ ساله}.

فازهای بیماری:
\begin{itemize}
\item[\en{a)}] کاتارال: علائم غیراختصاصی (تب خفیف) \rel \hl{۱ تا ۲ هفته}.
\item[\en{b)}] پاروکسیسمال: سرفه‌های شدید \rel \hl{۲ تا ۴ هفته}.

سرفه حین بازدم، استفراغ به دنبال سرفه.
\item[\en{c)}] کانوالسنت: بهبودی تدریجی \rel \hl{۶ تا ۸ هفته}.
\end{itemize}

\begin{itemize}
\item در نوزادان اولین علامت ممکن است \hl{آپنه} باشد.
\item در شیرخواران:
\begin{itemize}
\item آسیب مغزی ناشی از هایپوکسی.
\item عفونت باکتریال ثانویه.
\end{itemize}
\item در نوجوانان \rel برونشیت طول‌کشیده.
\end{itemize}

\subsection*{تشخیص بیماری}
\hl{کشت، \en{PCR}}.

\begin{itemize}
\item در ۸۰\% بیماران \hl{لنفوسیتوز} وجود دارد.
\en{WBC} \rel \hl{۲۰٬۰۰۰ تا ۵۰٬۰۰۰}.
\item در تصویربرداری:
\begin{itemize}
\item آتلکتازی سگمنتال می‌دهد.
\item انفیلتراسیون پری‌هیلار \rel شبیه پنومونی وایرال.
\end{itemize}
\end{itemize}

\clearpage
\sheet{۲}

\subsection*{درمان}
\begin{itemize}
\item \hl{ماکرولید} (آزیترو، کلاریترو، اریترو).
\begin{itemize}
\item \hl{به نوزادان آزیترو} می‌دهیم.
\item اریترو به نوزادان \rel استنوز پیلور.
\end{itemize}
\item درمان در فاز کاتارال \rel ظرف ۳ تا ۴ روز نازوفارنکس را پاک می‌کند.
\item درمان در فاز پاروکسیسمال \rel دورهٔ بیماری را تغییر نمی‌دهد ولی از انتشار به سایر کودکان جلوگیری می‌کند.
\item در بالای ۲ ماه \rel تریمتوپریم ـ سولفامتوکسازول می‌توانیم بدهیم.

{\small\textcolor{gray}{اصلاح دست‌نویس در همین سطر: \sout{سال}؛ «ماه» بالای آن نوشته شده است.}}
\end{itemize}

\subsection*{واکسن}
\begin{itemize}
\item در ۲، ۴، ۶، ۱۵ تا ۱۸ ماهگی، ۴ تا ۶ سالگی، ۱۱ تا ۱۲ سالگی.
\item در هر بارداری یک دوز بوستر در سه‌ماههٔ سوم.
\end{itemize}

\subsection*{پروفیلاکسی پس از مواجهه}
آزیترو ۵ روز / کلاریترو یا اریترو ۷ تا ۱۴ روز.

{\small\textcolor{gray}{\textbf{سطر خط‌خورده در منبع:}
\sout{اگر زیر ۷ سال} \rel \sout{یک دوز واکسن اگر واکسن نگرفته}.
[خوانش عبارت پایانی این سطر خط‌خورده نامطمئن است.]}}

\end{document}
